\begin{theorem}[Natural boundary of primitive--Dirichlet series: \texorpdfstring{$\Re(s)=0$}{Re(s)=0}]\label{thm:real-input-40-primedirichlet-dense-branch}
Let
\[
\mathcal{P}(s):=P(\lambda^{-s})=\sum_{n\ge 1}p_n\,\lambda^{-ns}\qquad(\Re(s)>1),
\]
where $P(z)=\sum_{n\ge 1}p_n z^n$ is the primitive generating function, $\lambda=\rho(M)=\varphi^2$, and define
\[
\widehat{\zeta}(s):=\zeta_M(\lambda^{-s}).
\]
Then:
\begin{enumerate}
  \item \textbf{(Analytic continuation to right half-plane)} For any $\sigma_0>0$, the series
  \[
  \sum_{k\ge 1}\frac{\mu(k)}{k}\log \widehat{\zeta}(ks)
  \]
  converges uniformly on any compact subset of
  \(
  \{s:\Re(s)\ge \sigma_0\}
  \)
  (avoiding the pole sequence below), thus defining an analytic continuation of $\mathcal{P}(s)$ on the open half-plane $\Re(s)>0$ (which coincides with the original Dirichlet series for $\Re(s)>1$).
  \item \textbf{(Dense singularity sequence and natural boundary)} For each odd prime $p$ and any $m\in\ZZ$, the point
  \[
  s_{p,m}:=\frac{1}{p}+\frac{2\pi i m}{p\log\lambda}
  \]
  is a logarithmic branch singularity of this analytic continuation. Moreover, the set $\{s_{p,m}\}_{p\ \mathrm{odd\ prime},\,m\in\ZZ}$ is dense on the line $\Re(s)=0$.
\end{enumerate}
Therefore \(\Re(s)=0\) is a \textbf{natural boundary} of \(\mathcal{P}(s)\): \(\mathcal{P}(s)\) cannot be analytically continued across this line in any neighborhood of any imaginary axis point.
\end{theorem}

\begin{proof}
\emph{Point 1 (Convergence and analyticity)}:
When $|z|$ is sufficiently small we have $\log\zeta_M(z)=O(z)$ (since $\zeta_M(0)=1$ and $\zeta_M$ is analytic in a neighborhood of $0$), hence for any $\sigma_0>0$, when $\Re(s)\ge \sigma_0$ we have
\(
|\log\widehat{\zeta}(ks)|=|\log\zeta_M(\lambda^{-ks})|=O(\lambda^{-k\sigma_0})
\)
decaying exponentially in $k$. Thus $\sum_{k\ge 1}|\mu(k)|\,|\log\widehat{\zeta}(ks)|/k$ converges uniformly on any compact set avoiding the pole sequence; term-by-term differentiability and uniform convergence show that the sum function is analytic on $\Re(s)>0$ (minus the pole sequence), yielding the analytic continuation of $\mathcal{P}(s)$.

\emph{Point 2 (Explicit singularities and density)}:
From the determinant factorization for this kernel
\[
\det(I-zM)=(1-z)^2(1+z)^3(1-3z+z^2)(1+z-z^2),
\]
we see $z_\star=\lambda^{-1}$ is a \textbf{simple pole} of $\zeta_M(z)$ (corresponding to the simple zero of the factor $1-3z+z^2$), so $\log\zeta_M(z)$ has a logarithmic branch singularity at $z=z_\star$.
When $s=s_{p,m}$ we have
\[
\lambda^{-ps}=\lambda^{-1}e^{-2\pi i m}=z_\star,
\]
hence $\log\widehat{\zeta}(ps)=\log\zeta_M(\lambda^{-ps})$ has a logarithmic branch singularity at $s=s_{p,m}$, thus the term
\(
\frac{\mu(p)}{p}\log\widehat{\zeta}(ps)
\)
is singular at this point. For any $k\neq p$, if $\log\widehat{\zeta}(ks)$ is singular at $s=s_{p,m}$, then we must have $\lambda^{-ks_{p,m}}=z_\star$, equivalent to $ks_{p,m}=1+\frac{2\pi i \ell}{\log\lambda}$ (for some $\ell\in\ZZ$). Comparing real parts gives $k/p=1$, so $k=p$. Therefore except for the $k=p$ term, all other terms are analytic in some neighborhood of $s_{p,m}$, and by the uniform convergence from Point 1 their sum is also analytic, so $\mathcal{P}(s)$ must be a logarithmic branch singularity at $s_{p,m}$.

Finally, to prove density: given any $t\in\RR$ and $\varepsilon>0$, take sufficiently large odd prime $p$, and let
\(
m:=\left\lfloor \frac{t\,p\log\lambda}{2\pi}\right\rceil
\)
be the nearest integer, then
\[
\left|\Im(s_{p,m})-t\right|
=\left|\frac{2\pi m}{p\log\lambda}-t\right|
\le \frac{\pi}{p\log\lambda}<\varepsilon
\]
holds for sufficiently large $p$, while $\Re(s_{p,m})=1/p<\varepsilon$. Thus $s_{p,m}\to it$, so $\{s_{p,m}\}$ is dense on the imaginary axis. Since any neighborhood crossing the imaginary axis contains these singularities, analytic continuation cannot cross, making \(\Re(s)=0\) a natural boundary.
\end{proof}

\begin{proposition}[Branch point inheritance (prime-zeta scaling law)]\label{prop:primedirichlet-branch-inherit}
Maintaining notation from Theorem \ref{thm:real-input-40-primedirichlet-dense-branch}, and assuming \(\Re(s_0)>0\).
If \(\widehat{\zeta}(s)\) has a pole at \(s=s_0\) (or equivalently \(\log\widehat{\zeta}(s)\) has a logarithmic branch singularity at \(s=s_0\)), then for any odd prime \(p\), the function \(\mathcal{P}(s)\) has a logarithmic branch singularity at
\[
s=\frac{s_0}{p}.
\]
\end{proposition}
\begin{proof}
From the prime-zeta form identity
\(
\mathcal{P}(s)=\sum_{k\ge 1}\frac{\mu(k)}{k}\log\widehat{\zeta}(ks)
\)
we know: at point \(s=s_0/p\), the term \(\frac{\mu(p)}{p}\log\widehat{\zeta}(ps)\) is a logarithmic branch singularity since \(ps=s_0\).
For any \(k\neq p\), if \(\log\widehat{\zeta}(ks)\) is singular at \(s=s_0/p\), then \(ks\) must lie on the pole sequence of \(\widehat{\zeta}\), in particular satisfying \(\Re(ks)=\Re(s_0)\).
Comparing real parts gives \(k\Re(s_0)/p=\Re(s_0)\), thus \(k=p\), contradiction.
Therefore except for the \(k=p\) term, all other terms are analytic in some neighborhood of \(s_0/p\); by uniform convergence from Point 1 of Theorem \ref{thm:real-input-40-primedirichlet-dense-branch}, the sum of analytic terms remains analytic, so \(\mathcal{P}(s)\) must be a logarithmic branch singularity at \(s=s_0/p\). This completes the proof.
\end{proof}

\begin{corollary}[Family of branch points at \texorpdfstring{$\Re(s)=\frac{1}{2p}$}{Re(s)=1/(2p)} induced by toy-RH critical line poles]\label{cor:real-input-40-primedirichlet-branch-halfp}
Maintaining notation from Theorem \ref{thm:real-input-40-primedirichlet-dense-branch}. For each odd prime \(p\) and any \(m\in\ZZ\), the point
\[
t_{p,m}:=\frac{1}{2p}+\frac{(2m+1)\pi i}{p\log\lambda}
\]
is a logarithmic branch singularity of the analytic continuation of \(\mathcal{P}(s)\) on \(\Re(s)>0\).
\end{corollary}
\begin{proof}
By Proposition \ref{prop:finite-rh-40}, \(\widehat{\zeta}(s)\) has a pole at
\(
s_0=\frac12+\frac{(2m+1)\pi i}{\log\lambda}
\)
with \(\Re(s_0)=1/2>0\). Taking \(t_{p,m}=s_0/p\), the conclusion follows immediately from Proposition \ref{prop:primedirichlet-branch-inherit}. This completes the proof.
\end{proof}

\begin{remark}[Universalization of natural boundary mechanism]\label{rem:primedirichlet-natural-boundary-universal}
The proof of Theorem \ref{thm:real-input-40-primedirichlet-dense-branch} uses only two structural facts:
\begin{enumerate}
  \item The pole of $\widehat{\zeta}(s)=\zeta_M(\lambda^{-s})$ at $s=1$ comes from the Perron root $\lambda$, so $\log\widehat{\zeta}(s)$ has a logarithmic branch singularity at this point;
  \item The prime-zeta form identity
  \[
  \mathcal{P}(s)=\sum_{k\ge 1}\frac{\mu(k)}{k}\log\widehat{\zeta}(ks)
  \]
  replicates this branch singularity proportionally via $s\mapsto s/k$ to $s_{p,m}=1/p+2\pi i m/(p\log\lambda)$ (odd primes $p$), with real part comparison via $k/p$ ruling out cancellation from different $k$.
\end{enumerate}
Therefore this "dense branch singularities \(\Rightarrow\) \(\Re(s)=0\) natural boundary" phenomenon does not depend on the specific factorization of this kernel;
as long as $M$ is a primitive nonnegative integer matrix (guaranteeing simple Perron root) and the primitive--Dirichlet series is defined in the same way, the completely isomorphic dense singularity mechanism will appear.
\end{remark}

\paragraph{From toy-RH to real RH research form: five auditable interface checklist}\label{par:finite-rh-rhstyle-toolkit}
To avoid repetition, this paragraph consolidates the interfaces already closed in this paper and aligned with RH form by research roadmap level into five items; each corresponds to citable proposition/theorem/appendix interfaces given in the text:
\begin{enumerate}
  \item \textbf{Spectrum \(\leftrightarrow\) critical line dictionary (geometric alignment)}:
  Via Dirichlet variable \(z=\lambda^{-s}\) one obtains pole vertical line real part identity \(\Re(s)=\log|\mu|/\log\lambda\), thus \(|\mu|=\sqrt{\lambda}\iff \Re(s)=1/2\) (see Proposition \ref{prop:s-plane-spectrum-poles} in Section \ref{sec:zeta-finite-part} and its Corollary \ref{cor:rh-stratification-splane}, consistent with the special case in Proposition \ref{prop:finite-rh-40} of this appendix).
  \item \textbf{Square-root error rigidity dichotomy (saturation \(\Rightarrow\) \(\Omega\)-oscillation)}:
  The finite-dimensional explicit formula template (Proposition \ref{prop:finite-dimensional-explicit-formula}) shows: once there exists a saturating spectral point \(|\mu|=\sqrt{\lambda}\), generally unimprovable square-root oscillation appears; on the real-input kernel this oscillation is produced by the unique saturating eigenvalue \(-\sqrt{\lambda}\) and is completely explicit (Proposition \ref{prop:finite-rh-40} and Corollary \ref{cor:real-input-40-chebyshev-explicit}).
  \item \textbf{Length \(\leftrightarrow\) scale prime theorem shape (\(x/\log x\) main term)}:
  The primitive main term \(p_n\sim \lambda^n/n\) under last-term dominance is equivalently rewritten as \(\pi(x)\sim \frac{\lambda\log\lambda}{\lambda-1}\frac{x}{\log x}\); constant factors and derivation details see paragraph \ref{par:finite-rh-primitive-xlogx}.
  \item \textbf{Limitization of toy-GRH (profinite/Haar equidistribution)}:
  Chebotarev estimates for finite modular axis bundles under kernel-version GRH (square-root bound) can be lifted to the profinite inverse limit axis: for any cylinder set giving Haar weight main term with square-root error (Theorem \ref{thm:pom-profinite-axis}).
  \item \textbf{Natural boundary of prime Dirichlet series ("prime layer direct series" non-continuation constraint)}:
  The primitive--Dirichlet series \(\mathcal{P}(s)\) exhibits dense branch singularities on \(\Re(s)=0\), so the imaginary axis is a natural boundary (Theorem \ref{thm:real-input-40-primedirichlet-dense-branch}), and this mechanism is universal in the general case of primitive nonnegative integer kernels (Remark \ref{rem:primedirichlet-natural-boundary-universal}).
\end{enumerate}
These five interface items provide a more RH-research-aligned cumulative roadmap: first increase critical angle diversity on the spectral side via lift/twist (Remark \ref{rem:finite-rh-angle-entropy}), while using Chebotarev/profinite axis to lift congruence equidistribution to limit caliber; and under the natural boundary constraint of prime-zeta, prioritize organizing controllable superposition of error and oscillation terms via \(\zeta\) (or its logarithmic derivative/explicit formula) rather than direct primitive series.

\paragraph{From finite kernel template to Ihara/Selberg/Hilbert--Pólya roadmap: spectrum \texorpdfstring{$\Rightarrow$}{=>} functional equation, critical line structure, and Ramanujan bounds}
This paper's "finite kernel \(\zeta=\det(I-zT)^{-1}\)" prototype structurally aligns directly with three mature spectral roadmaps; the differences lie only in "the category where the operator resides (finite matrix vs.\ infinite-dimensional kernel operator)" and "arithmetic origin of twisted indexing (character/covering vs.\ Hecke/prime)":
\begin{itemize}
  \item \textbf{Ihara/graph \(\zeta\) (Ramanujan \(\Leftrightarrow\) RH on graphs)}:
  For regular graphs, the "critical circle/critical line" condition of Ihara \(\zeta\) and adjacency spectrum satisfying Ramanujan bounds are equivalent paradigms (see \cite{TerrasZetaGraphs,StarkTerras1996GraphCoverings}). This belongs to the same mother form as this paper's kernel RH/GRH ("nontrivial spectral radius \(\le \lambda^{1/2}\)"): critical geometry of \(\zeta\) is controlled by spectral radius bounds, while covering/character twist organizes congruence decomposition into Artin-type factors (see unified caliber for Dirichlet stratification and twisted spectrum in this paper).
  \item \textbf{Selberg/Ruelle (\(\zeta\)'s Fredholm determinant and spectral interpretation)}:
  In hyperbolic geodesic flows and general Axiom A systems, Ruelle/Selberg type \(\zeta\) can be formulated as (generalized) Fredholm determinants of certain transfer operator families, whose zeros/poles correspond to spectral data (resonances/eigenvalues) of operators (see \cite{Ruelle1976ZetaExpanding,ParryPollicott1990Zeta,Fried1986RuelleSelberg}). Thus, if we want "imaginary part structure on critical line" to no longer degenerate to finite lattice, the most natural upgrade is to replace this paper's finite kernel \(T\) with a controllable family of kernel operators \(\mathcal{L}_u\), and inherit this paper's \(\zeta\)-interface via Fredholm determinant \(\det(I-z\mathcal{L}_u)\).
  \item \textbf{Hilbert--Pólya/Connes (explicit formula \(\leftrightarrow\) trace formula spectralization)}:
  One spectral interpretation roadmap for RH is to write the explicit formula as a trace formula, viewing critical zeros as some kind of "spectrum/resonance" (see \cite{Connes1999TraceFormulaNCG}). From this paper's perspective, this corresponds to lifting the "trace/primitive commutative diagram" already closed in the text to an operator category with arithmetic symmetry: organization of twisted channels no longer comes only from finite group characters, but is provided by operator families indexed by Hecke/primes.
\end{itemize}
Upon these three alignments, this paper's "two-variable self-dual—completion" structure (functional equation \(s\mapsto 1-s\) in Appendix \ref{app:kernel-functional-equation} and Hardy real-value) can serve as hard constraint for upgrade roadmap: first maintain self-dual/completion to lock functional equation and critical line geometry, then extend finite-dimensional spectral radius bound to infinite-dimensional Ramanujan-type spectral bound (kernel-version GRH), and use Artin/character decomposition and profinite axis to lift congruence equidistribution to limit caliber.

\paragraph{Alignment with analytic number theory zero-density techniques}
In \(\zeta\) theory, one cumulative roadmap toward RH is to improve large-value control for Dirichlet polynomials, thereby improving zero-density bounds \(N(\sigma,T)\) and feeding back to prime error term thresholds.
For example, Guth--Maynard give new large-value estimates and derive improved zero-density upper bounds and short interval prime counting thresholds (\(N(\sigma,T)\le T^{30(1-\sigma)/13+o(1)}\) and related corollaries, see \cite{GuthMaynard2024LargeValues}).

From the "alignment caliber" perspective, the role of zero-density techniques is not to directly derive "all zeros are on the critical line," but to thin the set of zeros deviating from the critical line, thereby controlling oscillating error terms in the explicit formula to stronger integrable/averageable forms.
This corresponds to the structure in this paper's finite-dimensional prototype where "worst twisted spectral radius ratio $\Lambda/\lambda_1$ determines error exponent": when square-root bound (kernel-version GRH) cannot be obtained in one step, one can still improve the error-controllable range by controlling appearance frequency/magnitude of "bad twists."

From this paper's "finite-dimensional \(\zeta\)" perspective, the above chain can be understood as a continuous-spectrum version "spectrum—pole—error triangle": sparsity of zero (or pole) distribution corresponds to controllable superposition of oscillating terms in the explicit formula, while the achievable error term exponent is controlled by available large-value bounds/moment mother functions. The advantage provided by the finite-dimensional model is: this triangle can be compressed completely auditability to finite spectral data (Proposition \ref{prop:finite-rh-triangle-dict} and Corollary \ref{cor:finite-rh-explicit-formula}), thus providing a computable prototype for "writing error exponent as verifiable spectral invariant."

\paragraph{Quantifiable boundaries for real \texorpdfstring{$\zeta$}{zeta}: hard inputs needed for alignment and roadmap design}
To interface external progress in RH direction into this paper's spectral language, here we list several hard inputs still directly callable in 2024--2025 (all verifiable explicit constants/thresholds):
\begin{itemize}
  \item \textbf{Verification height (partial RH):} Rigorously verified up to height \(3\cdot 10^{12}\), i.e., all nontrivial zeros satisfying \(0<\Im(\rho)\le 3\cdot 10^{12}\) lie on the critical line \(\Re(\rho)=\tfrac12\) (\cite{PlattTrudgian2020RH3e12}).
  \item \textbf{de Bruijn--Newman constant:} Rodgers--Tao proved \(\Lambda\ge 0\) (\cite{RodgersTao2018NewmanNonnegative}), while RH is equivalent to \(\Lambda\le 0\), hence RH is equivalent to \(\Lambda=0\). Meanwhile Polymath15 gives unconditional upper bound \(\Lambda\le 0.22\), with conditional improvements dependent on further numerical verification (\cite{Polymath2019NewmanUpper}).
  \item \textbf{Proportion of zeros on critical line:} Research shows over \(5/12\) of zeros lie on the critical line (\cite{PrattRoblesZaharescuZeindler2018FiveTwelfths}). Such "proportion lifting" can be viewed as one cumulative roadmap for thinning the set of zeros deviating from critical line.
  \item \textbf{Explicit zero counting error:} Bellotti--Wong improved explicit error bounds for \(N(T)\), giving specific constants (\cite{BellottiWong2024ImprovedNT}), directly usable for Turing method/numerical verification and explicit applications.
  \item \textbf{Improvement of Turing method key quantities:} Amberger gives improved bounds for \(|S_1(t_2)-S_1(t_1)|\), thereby improving block length thresholds required by Turing method (\cite{Amberger2025S1}).
  \item \textbf{Explicit control at \(\Re(s)\approx 1\):} Cully-Hugill--Leong and Leong respectively give explicit upper bounds for \(\zeta'(s)/\zeta(s)\) and \(1/\zeta(s)\) when \(\sigma\) approaches \(1\) (\cite{CullyHugillLeong2023Zeta1Line,Leong2024LogDerivRecip}), usable as engineering inputs for zero-free regions and error term constants.
\end{itemize}

\begin{definition}[Kernel-version Newman threshold (verifiable parameter)]\label{def:kernel-newman-threshold}
Given a parametric kernel family \(M(u)\) (e.g., generated by potential function/temperature gate \(\mathrm{PROJ}_u\) or character twist; see Definitions \ref{def:pom-proj-temp} and \ref{def:pom-proj-chi}), denote main spectral radius as \(\lambda_1(u)\), worst nontrivial twisted spectral radius as \(\Lambda(u)\).
Define kernel-version Newman threshold as
\[
\beta_\ast:=\inf\{\beta\in\RR_{\ge 0}:\ \Lambda(e^{-\beta})\le \lambda_1(e^{-\beta})^{1/2}\}.
\]
When \(\beta_\ast=0\), the origin slice (\(u=1\)) already satisfies kernel-version GRH square-root bound; when \(\beta_\ast>0\), one needs to deform toward zero temperature (\(\beta\) increases, \(u=e^{-\beta}\) decreases) to enter square-root regime.
\end{definition}

\begin{remark}[Reproducible output: kernel-version Newman threshold for pure collision axis (finite modulus proxy)]\label{rem:kernel-newman-threshold-repro}
Taking the three-axis twist family of the real-input 40-state kernel as example, fix \((q,r)=(1,1)\), and on the third axis take pure collision observable \(N_2\) (i.e., \(\xi=\mathbf{1}_{d=2}\)), for each odd prime \(p\) examine composite twist
\(
u=\exp(-\beta)\,\omega_p^k
\)
(\(k\not\equiv 0\!\!\pmod p\)). Define
\[
\lambda_1(\beta):=\rho(M(\exp(-\beta))),\qquad
\Lambda_p(\beta):=\max_{k\neq 0}\rho\bigl(M(\exp(-\beta)\omega_p^k)\bigr),
\]
and let \(\beta_\ast(p)\) be the minimal \(\beta\) (if exists) satisfying \(\Lambda_p(\beta)\le \lambda_1(\beta)^{1/2}\).
Table \ref{tab:real_input_40_kernel_newman_threshold} gives numerical scan output for several \(p\): if no \(\beta\) satisfying square-root bound is observed in scan interval, give lower bound \(\beta_\ast(p)>\beta_{\max}\), while also listing minimum value of \(\Lambda_p/\sqrt{\lambda_1}\) on scan grid and its occurrence location (script \texttt{scripts/exp\_real\_input\_40\_kernel\_newman\_threshold.py} one-click generation):
\input{sections/generated/tab_real_input_40_kernel_newman_threshold_en}
\end{remark}

\begin{remark}[From toy-RH to "angle entropy": an operational upgrade direction]\label{rem:finite-rh-angle-entropy}
Proposition \ref{prop:finite-rh-triangle-dict} shows: the "critical line position" of finite-dimensional toy-RH depends only on modulus constraint of normalized spectrum (\( |\mu|\le \sqrt\lambda \)).
To simulate closer to \(\zeta\)'s "zero density/statistical form" in finite models, one must let eigenvalues at critical modulus (\(|\mu|=\sqrt\lambda\)) carry many different angles, thereby extending single lattice oscillation to multi-frequency oscillation superposition (structural extension of Corollary \ref{cor:finite-rh-explicit-formula}).
Under this paper's compilation caliber, this goal can be achieved by increasing lift/character axis bundles (e.g., profinite congruence axis or higher-dimensional abelian twist families), and "entropy/diversity of unit circle spectrum angle distribution" can be used as reproducible experimental signature.
A more direct "multi-frequency fluctuation" readout is primitive-sharp \(\mathsf{E}_n^\sharp\) stripped of prime-powers (Definition \ref{def:primitive-sharp}): it strips Perron's prime-power interference terms, compressing square-root error to critical modulus phase sum; when phases are rationally independent it also gives limit distribution of Haar pushforward (Remark \ref{rem:primitive-sharp-limit-law}).
\end{remark}

\begin{remark}[Auditable certificate for square-root bound: Lyapunov/inner product construction]\label{rem:finite-rh-lyapunov}
Kernel-version (G)RH square-root bound is a spectral radius inequality. Besides direct diagonalization/eigenvalue computation, one can also use "inner product construction" to give auditable certificate:
For any matrix $A$ and threshold $r$, $\rho(A)<r$ is equivalent to existence of positive definite matrix $H\succ 0$ such that $A^\ast H A\prec r^2 H$ (Proposition \ref{prop:lyapunov-spectral-radius-certificate}).
Thus on twisted component $A=M_{K,\rho}(u)$, if one can explicitly give certificate matrix $H_{\rho,u}$, the square-root bound can be written as a verifiable inequality, independent of numerical eigenvalue sorting robustness.
\par
More importantly: the inequality
\(
A^\ast H A\prec r^2 H
\)
is a linear matrix inequality (LMI) for unknown \(H\), hence can be directly rewritten as semidefinite feasibility problem (SDP).
This provides an "auditable + automatically searchable" research loop:
\begin{enumerate}
  \item Fix a lift/labeling skeleton (e.g., permutation lift; see Artin factorization framework in Proposition \ref{prop:kernel-artin-factorization}) and treat its free parameters (generators, permutations, cocycles) as variables to optimize;
  \item For each irreducible channel \(\rho\neq\mathbf{1}\), construct corresponding twisted block matrix \(M_{K,\rho}(u)\);
  \item For each channel solve SDP seeking \(H_{\rho,u}\succ 0\) such that
  \[
  M_{K,\rho}(u)^\ast H_{\rho,u} M_{K,\rho}(u)\prec \lambda_1(u)\,H_{\rho,u}
  \qquad(\text{i.e., }r=\sqrt{\lambda_1(u)}),
  \]
  if all nontrivial channels are feasible, one obtains an \textbf{auditable certificate object} for kernel-version GRH (the \(H\) matrix itself can be independently verified).
\end{enumerate}
This "certificate-feasibility" perspective rewrites the assertion "zeros/poles on critical line" as feasibility verification of finitely many SDPs, thus particularly suitable for systematic search and structural induction: first use computation to find lifts satisfying (or approximating) square-root bound, then abstract general theorems from certificate shapes.
\end{remark}

\paragraph{Primitive orbit counting curve and constant factor of \texorpdfstring{$x/\log x$}{x/log x}}\label{par:finite-rh-primitive-xlogx}
Let
\[
\Pi(N):=\sum_{n\le N}p_n
\]
be the cumulative primitive orbit count by length (counting by orbit). From $p_n\sim \lambda^n/n$ "last-term dominance + geometric tail" one obtains
\[
\Pi(N)\sim \frac{\lambda}{\lambda-1}\cdot \frac{\lambda^N}{N}.
\]
On this kernel $\lambda=\varphi^2$ and $\lambda-1=\varphi$, so
\[
\Pi(N)\sim \varphi\cdot\frac{\varphi^{2N}}{N}=\frac{\varphi^{2N+1}}{N}.
\]
If one defines "orbit norm" $N(\tau):=\lambda^{|\tau|}$ and lets
\[
\pi(x):=\#\{\tau\ \mathrm{primitive}:\ N(\tau)\le x\},
\]
then $\pi(\lambda^N)=\Pi(N)$, thus
\[
\pi(x)\sim \frac{\lambda\log\lambda}{\lambda-1}\cdot \frac{x}{\log x}.
\]
For $\lambda=\varphi^2$, this constant factor simplifies to
\[
\frac{\lambda\log\lambda}{\lambda-1}
=\frac{\varphi^2\log(\varphi^2)}{\varphi}
=2\varphi\log\varphi.
\]

Define orbit Mertens sum
\[
\mathcal{M}(N):=\sum_{n\le N}\frac{p_n}{\lambda^n}
=\log N+\mathsf{M}_{\varphi^2}+o(1),
\]
where the constant term can be written as
\[
\mathsf{M}_{\varphi^2}
=\gamma+\sum_{n\ge 1}\left(\frac{p_n}{\lambda^n}-\frac{1}{n}\right),
\]
$\gamma$ being Euler--Mascheroni constant. Numerically
\[
\sum_{n\ge 1}\left(\frac{p_n}{\lambda^n}-\frac{1}{n}\right)\approx -0.254312661901716,
\qquad
\mathsf{M}_{\varphi^2}\approx 0.322903002999817.
\]
Also denote $P(z)=\sum_{n\ge 1}p_n z^n$, then as $z\to z_\star=\lambda^{-1}$
\[
P(z)=-\log(1-\lambda z)+\log\mathfrak{M}+o(1),
\]
where
\[
\log\mathfrak{M}\approx -0.254312661901716,\qquad
\mathfrak{M}\approx 0.7754493106833835.
\]
From factorization one obtains main pole $z_\star=\lambda^{-1}=\varphi^{-2}$, and
\[
\zeta_M(z)=\frac{C}{1-\lambda z}+O(1),\qquad
C=\lim_{z\to z_\star}(1-\lambda z)\,\zeta_M(z)=\frac{47+21\sqrt5}{100}\ (\approx 0.9395742753).
\]

\begin{proposition}[Closed form of main pole residue constant (auditable)]\label{prop:real-input-40-residue-constant}
Continuing notation \(\lambda=\rho(M)=\varphi^2\) and \(z_\star=\lambda^{-1}\) from above, and defining main pole residue constant
\[
C_{\mathrm{real}}:=\lim_{z\to z_\star}(1-\lambda z)\,\zeta_M(z).
\]
Then \(C_{\mathrm{real}}\) has closed form
\[
\boxed{\;
C_{\mathrm{real}}=\frac{47+21\sqrt5}{100}=\frac{13+21\varphi}{50}\; }.
\]
\end{proposition}

\begin{corollary}[Dynamical system version Chebyshev explicit formula: main term \(\boldsymbol{+}\) square-root oscillation term]\label{cor:real-input-40-chebyshev-explicit}
Denote
\[
\Lambda(n):=a_n=\Tr(M^n),
\qquad
\Psi(N):=\sum_{n\le N}\Lambda(n).
\]
Then from spectral closed form
\(
a_n=\varphi^{2n}+(-\varphi)^n+2+3(-1)^n+\varphi^{-2n}+\varphi^{-n}
\)
summing geometric series term-by-term gives completely exact identity (no \(O(\cdot)\) terms):
\[
\boxed{\;
\begin{aligned}
\Psi(N)
={}&\frac{\lambda^{N+1}-\lambda}{\lambda-1}
+\frac{1-\lambda^{-N}}{\lambda-1}
+\frac{(-\varphi)^N-1}{\varphi}
+2N+\frac{3}{2}\bigl((-1)^N-1\bigr)
+\bigl(\varphi-\varphi^{1-N}\bigr),
\end{aligned}}
\]
where \(\lambda=\varphi^2\).
Hence the "RH-sharp" square-root scale oscillation term has explicit main amplitude:
\[
\frac{(-\varphi)^N}{\varphi}=(-1)^N\frac{1}{\varphi}\,\lambda^{N/2},
\]
whose magnitude is exactly square root of main term scale, with strictly alternating sign.
\end{corollary}
\begin{proof}
From \(\zeta_M(z)=1/\Delta(z)\) and \(z_\star\) being simple zero of \(\Delta\), we get
\(
\zeta_M(z)=\frac{1}{-\Delta'(z_\star)}\cdot\frac{1}{z-z_\star}+O(1)
\).
Also \(1-\lambda z=-\lambda(z-z_\star)+O((z-z_\star)^2)\), so
\(
C_{\mathrm{real}}=\frac{\lambda}{\Delta'(z_\star)}=\frac{1}{-z_\star\Delta'(z_\star)}
\).
Substituting factor decomposition of \(\Delta(z)=\det(I-zM)\) and simplifying gives the displayed algebraic closed form. This completes the proof.
\end{proof}
Further by Witt decomposition/ Möbius--log inversion (same caliber as \S\ref{subsec:zeta-online-kernel}) one obtains an Abel/finite-part closed form directly read from $\zeta_M$ main pole:
\[
\log\mathfrak{M}
=\log C+\sum_{m\ge 2}\frac{\mu(m)}{m}\log\zeta_M(z_\star^{\,m}),
\qquad z_\star=\lambda^{-1}.
\]
The right side substitutes rational function $\zeta_M$ into algebraic numbers $z_\star^{\,m}$; since $z_\star<1$, this series converges exponentially.

\paragraph{Explicit Möbius--$\zeta$ series and logarithmic product (using only factorization of \texorpdfstring{$\det(I-zM)$}{det})}
From
\[
\zeta_M(z)=\frac{1}{\Delta(z)},\qquad
\Delta(z)=\det(I-zM)=(1-z)^2(1+z)^3(1-3z+z^2)(1+z-z^2),
\]
and $z_\star=\lambda^{-1}\in(0,1)$, the above can be equivalently written in form containing only $\Delta$
\[
\log\mathfrak{M}
\;=\;
\log C-\sum_{m\ge 2}\frac{\mu(m)}{m}\log \Delta(z_\star^{\,m}).
\]
Further expanding $\Delta$ factors term-by-term gives completely "positive real logarithm" closed form:
\[
\log\mathfrak{M}
=\log C
-\sum_{m\ge 2}\frac{\mu(m)}{m}\Bigl[
2\log(1-z_\star^{\,m})+3\log(1+z_\star^{\,m})+\log(1-3z_\star^{\,m}+z_\star^{\,2m})+\log(1+z_\star^{\,m}-z_\star^{\,2m})
\Bigr].
\]
This formula is extremely numerically stable: since $z_\star^{\,m}$ tends to zero exponentially, and each term only requires taking logarithm of positive numbers.
Equivalently, can be written as Euler--Möbius product (compressing "finite part constant" completely to a comparable scalar):
\[
\mathfrak{M}
=C\prod_{m\ge 2}\zeta_M(z_\star^{\,m})^{\mu(m)/m}
=C\prod_{m\ge 2}\Delta(z_\star^{\,m})^{-\mu(m)/m}.
\]
\begin{lemma}[Möbius collapse: explicit summation of trivial factors]\label{lem:mobius-collapse}
For $|z|<1$ there hold identities
\[
\sum_{m\ge 1}\frac{\mu(m)}{m}\log(1-z^m)=-z,\qquad
\sum_{m\ge 1}\frac{\mu(m)}{m}\log(1+z^m)=z-z^2.
\]
Therefore
\[
\sum_{m\ge 1}\frac{\mu(m)}{m}\log\Bigl((1-z^m)^{-a}(1+z^m)^{-b}\Bigr)=(a-b)z+bz^2.
\]
\end{lemma}

\begin{proof}
The first identity can be directly verified by power series expansion: when $|z|<1$
\[
\sum_{m\ge 1}\frac{\mu(m)}{m}\log(1-z^m)
=-\sum_{m\ge 1}\frac{\mu(m)}{m}\sum_{k\ge 1}\frac{z^{mk}}{k}
=-\sum_{n\ge 1}\frac{z^n}{n}\sum_{m\mid n}\mu(m)
=-z.
\]
Here we used the identity $\sum_{m\mid n}\mu(m)=\mathbf{1}_{\{n=1\}}$.

The second identity follows from
\(
\log(1+z^m)=\log(1-z^{2m})-\log(1-z^m)
\)
reducing to the first: let
\(
f(z):=\sum_{m\ge 1}\frac{\mu(m)}{m}\log(1-z^m)
\),
then
\[
\sum_{m\ge 1}\frac{\mu(m)}{m}\log(1+z^m)=f(z^2)-f(z)=(-z^2)-(-z)=z-z^2.
\]
The last formula follows directly from linear combination
\(
-a\log(1-z^m)-b\log(1+z^m)
\)
and the first two. This completes the proof.
\end{proof}

\begin{lemma}[Explicit summation starting from \texorpdfstring{$m\ge 2$}{m>=2}]\label{lem:mobius-collapse-mge2}
For $|z|<1$ there hold identities
\[
\sum_{m\ge 2}\frac{\mu(m)}{m}\log(1-z^m)=-z-\log(1-z),
\]
\[
\sum_{m\ge 2}\frac{\mu(m)}{m}\log(1+z^m)=z-z^2-\log(1+z).
\]
\end{lemma}
\begin{proof}
From Lemma \ref{lem:mobius-collapse}'s two identities, directly writing
\(
\sum_{m\ge 2}=\sum_{m\ge 1}-\text{(}m=1\text{ term)}
\)
gives the result. This completes the proof.
\end{proof}

\begin{corollary}[General form: \texorpdfstring{$m\ge 2$}{m>=2} Möbius--$\log$ closed form for \texorpdfstring{$(1-z)^a(1+z)^b$}{(1-z)^a(1+z)^b}]\label{cor:mobius-collapse-mge2-ab}
For $|z|<1$, let
\[
S_{a,b}(z):=\sum_{m\ge 2}\frac{\mu(m)}{m}\log\Bigl((1-z^m)^{-a}(1+z^m)^{-b}\Bigr).
\]
Then there holds strict identity
\[
\boxed{\;
S_{a,b}(z)=(a-b)z+bz^2+a\log(1-z)+b\log(1+z)\; }.
\]
\end{corollary}

\begin{proof}
By definition
\(
S_{a,b}(z)=-a\sum_{m\ge 2}\frac{\mu(m)}{m}\log(1-z^m)-b\sum_{m\ge 2}\frac{\mu(m)}{m}\log(1+z^m)
\).
Substituting Lemma \ref{lem:mobius-collapse-mge2}'s two closed forms and rearranging gives the result. This completes the proof.
\end{proof}

\begin{remark}[Nontrivial part after stripping trivial factors]
Applying Lemma \ref{lem:mobius-collapse} to $(1-z)^2(1+z)^3$, its Möbius--$\log$ summation contribution can be precisely compressed to
\[
(-z_\star+3z_\star^2),
\]
where $-z_\star$ comes from linear merged term of $(1-z)^2$ and $(1+z)^3$, while $+3z_\star^2$ is quadratic correction term of $(1+z)^3$ in correction identity.
At this kernel's main pole $z_\star=\varphi^{-2}$ this further simplifies to complete algebraic value
\[
-z_\star+3z_\star^2
=-\varphi^{-2}+3\varphi^{-4}
=9-4\sqrt5
\approx 0.055728090000840.
\]
Therefore "stripping trivial factors" not only produces linear term, but also brings an explicitly computable quadratic term; this does not change main pole position, but produces observable constant-level correction in compressed expression of Abel finite part constant.
\end{remark}

\begin{proposition}[Primitive contribution of trivial factors is finite polynomial]\label{prop:triv-factor-primitive-polynomial}
Define
\[
\zeta_{\mathrm{triv}}(z):=\frac{1}{(1-z)^2(1+z)^3},
\qquad
P_{\mathrm{triv}}(z):=\sum_{m\ge 1}\frac{\mu(m)}{m}\log\zeta_{\mathrm{triv}}(z^m).
\]
Then for \(|z|<1\) there holds exact identity
\[
\boxed{\ P_{\mathrm{triv}}(z)=-z+3z^2\ }.
\]
Hence trivial factors' influence on primitive orbit generating function \(P(z)=\sum_{n\ge 1}p_n z^n\) occurs only at \(n=1,2\):
\[
p^{\mathrm{triv}}_1=-1,\qquad p^{\mathrm{triv}}_2=3,\qquad p^{\mathrm{triv}}_n=0\ (n\ge 3),
\]
where "negative sign" only represents algebraic splitting under Witt/Möbius--$\log$ transform (not corresponding to negative counting of real orbits).
\end{proposition}
\begin{proof}
By Lemma \ref{lem:mobius-collapse},
\[
P_{\mathrm{triv}}(z)
=-2\sum_{m\ge 1}\frac{\mu(m)}{m}\log(1-z^m)-3\sum_{m\ge 1}\frac{\mu(m)}{m}\log(1+z^m)
=-2(-z)-3(z-z^2)=-z+3z^2.
\]
Comparing coefficients gives the conclusion for \(p^{\mathrm{triv}}_n\). This completes the proof.
\end{proof}

\begin{proposition}[Closed forms of two Möbius--$\log$ series: two-kernel compression of \texorpdfstring{$\log\mathfrak{M}$}{logM}]\label{prop:real-input-40-mertens-two-series}
On the real-input $40$ state kernel,
\(
\Delta(z)=(1-z)^2(1+z)^3(1-3z+z^2)(1+z-z^2)
\)
with \(z_\star=\lambda^{-1}=\varphi^{-2}\in(0,1)\), \(C=\frac{47+21\sqrt5}{100}\).
Define
\[
\log K_0
:=\log C-z_\star+3z_\star^2+2\log(1-z_\star)+3\log(1+z_\star)
=\log\!\left(\frac{5+2\sqrt5}{10}\right)+(9-4\sqrt5).
\]
Then the Abel finite part constant satisfies strict closed form
\[
\boxed{\;
\log\mathfrak{M}
=
\log K_0
-\sum_{m\ge 2}\frac{\mu(m)}{m}
\log\!\Bigl(
(1-3z_\star^{\,m}+z_\star^{\,2m})(1+z_\star^{\,m}-z_\star^{\,2m})
\Bigr)\; }.
\]
\end{proposition}
\begin{proof}
From the above formula
\(
\log\mathfrak{M}
=\log C-\sum_{m\ge 2}\frac{\mu(m)}{m}\log \Delta(z_\star^{\,m})
\)
expanding \(\Delta\)'s four factors term-by-term gives
\[
\log\mathfrak{M}
=\log C
-2\sum_{m\ge 2}\frac{\mu(m)}{m}\log(1-z_\star^{\,m})
-3\sum_{m\ge 2}\frac{\mu(m)}{m}\log(1+z_\star^{\,m})+(\cdots),
\]
where \((\cdots)\) denotes Möbius--$\log$ summation terms for two quadratic factors \(1-3z+z^2\) and \(1+z-z^2\).
For the first two sums, applying Lemma \ref{lem:mobius-collapse-mge2} immediately gives
\[
-2\sum_{m\ge 2}\frac{\mu(m)}{m}\log(1-z_\star^{\,m})
-3\sum_{m\ge 2}\frac{\mu(m)}{m}\log(1+z_\star^{\,m})
=-z_\star+3z_\star^2+2\log(1-z_\star)+3\log(1+z_\star),
\]
thus obtaining the displayed two-kernel compression formula. This completes the proof.
\end{proof}

\begin{remark}[Reproducible numerics: \texorpdfstring{$\log K_0,S_3,S_{\pm},\log\mathfrak{M}$}{logK0,S3,Spm,logM}]\label{rem:real-input-40-mertens-two-series-numeric}
Table \ref{tab:real_input_40_abel_mertens_two_series_split} gives high-precision numerics for two-kernel compression decomposition. We see
\(\log K_0\) has magnitude around \(10^{-3}\), while main contribution to \(\log\mathfrak{M}\) comes from sum of two Möbius--$\log$ series
corresponding to two quadratic factors \((1-3z+z^2)\) and \((1+z-z^2)\).
\input{sections/generated/tab_real_input_40_abel_mertens_two_series_split_en}
\end{remark}

\begin{theorem}[Complete closed-form stripping of trivial factors and \texorpdfstring{$\mathfrak{M}_{\mathrm{carry}}$}{M_carry} (quadratic carry-constant)]\label{thm:real-input-40-carry-constant}
Let \(z_\star=\varphi^{-2}\), \(C=\frac{47+21\sqrt5}{100}\).
Define
\[
\mathfrak{M}_{\mathrm{carry}}
:=\exp\!\left(
-\sum_{m\ge 2}\frac{\mu(m)}{m}
\Bigl[
\log(1-3z_\star^{\,m}+z_\star^{\,2m})
\;+\;
\log(1+z_\star^{\,m}-z_\star^{\,2m})
\Bigr]\right).
\]
Then the Abel finite part constant satisfies multiplicative decomposition
\[
\boxed{\;
\mathfrak{M}
\;=\;
\underbrace{C\cdot z_\star(1+z_\star)^3\cdot e^{z_\star^3}}_{\text{completely closed-form trivial factor part}}
\cdot
\underbrace{\mathfrak{M}_{\mathrm{carry}}}_{\text{carry-constant contributed only by two quadratic factors}}\; }.
\]
\end{theorem}

\begin{proof}
By Proposition \ref{prop:real-input-40-mertens-two-series},
\(
\log\mathfrak{M}
=\log K_0-(S_3+S_\pm)
\),
where
\(
\log K_0=\log C-z_\star+3z_\star^2+2\log(1-z_\star)+3\log(1+z_\star)
\).
Note that \(z_\star=\varphi^{-2}\) satisfies \(1-z_\star=\varphi^{-1}\), thus \((1-z_\star)^2=z_\star\); and from \(z_\star^2-3z_\star+1=0\) we get
\(
-z_\star+3z_\star^2=z_\star^3
\).
Hence
\[
K_0=C\cdot (1-z_\star)^2(1+z_\star)^3\cdot e^{-z_\star+3z_\star^2}
=C\cdot z_\star(1+z_\star)^3\cdot e^{z_\star^3}.
\]
Substituting \(\mathfrak{M}_{\mathrm{carry}}:=\exp(-(S_3+S_\pm))\) back gives the multiplicative decomposition. This completes the proof.
\end{proof}

\begin{remark}[Reproducible numerics and certified error bounds (carry-constant splitting)]\label{rem:real-input-40-carry-constant-certified}
Script \texttt{scripts/exp\_real\_input\_40\_carry\_constant\_certified.py} directly computes by Theorem \ref{thm:real-input-40-carry-constant}
\(\log\mathfrak{M}\), \(\log\mathfrak{M}_{\mathrm{carry}}\)
and explicit tail bound
\(
|\mathrm{tail}(N)|\le \frac{5 z_\star^{N+1}}{(N+1)(1-z_\star)}
\), outputting a directly embeddable formula block:
\input{sections/generated/eq_real_input_40_carry_constant_certified_en}
\end{remark}

\begin{remark}[Artin decomposition at constant level: three-factor multiplicative factorization of \texorpdfstring{$\log\mathfrak{M}$}{logM}]\label{rem:real-input-40-logM-artin-factorization}
From
\(
\Delta(z)=(1-z)^2(1+z)^3(1-3z+z^2)(1+z-z^2)
\)
we can decompose \(\zeta_M(z)=1/\Delta(z)\) into three factors
\[
\zeta_M(z)=\zeta_{\mathrm{triv}}(z)\cdot \zeta_{\mathrm{even}}(z)\cdot \zeta_{\chi}(z),
\]
where
\[
\zeta_{\mathrm{triv}}(z)=\frac{1}{(1-z)^2(1+z)^3},\qquad
\zeta_{\mathrm{even}}(z)=\frac{1}{1-3z+z^2},\qquad
\zeta_{\chi}(z)=\frac{1}{1+z-z^2}.
\]
At main pole \(z_\star=\lambda^{-1}=\varphi^{-2}\), only \(\zeta_{\mathrm{even}}\) has a pole, the other two factors are analytic and nonzero; hence residue and Abel finite part constant both admit multiplicative decomposition. Specifically, denoting
\[
P_{\mathrm{triv}}(z):=\sum_{m\ge 1}\frac{\mu(m)}{m}\log\zeta_{\mathrm{triv}}(z^m),\qquad
K_{-F}:=\sum_{m\ge 1}\frac{\mu(m)}{m}\log\zeta_{\chi}(z_\star^{\,m}),
\]
there exists \(\log\mathfrak{M}_{\mathrm{even}}\) such that
\[
\boxed{\ \log\mathfrak{M}=\log\mathfrak{M}_{\mathrm{even}}+K_{-F}+P_{\mathrm{triv}}(z_\star)\ }.
\]
Here \(P_{\mathrm{triv}}(z)\) is given by Proposition \ref{prop:triv-factor-primitive-polynomial} as finite polynomial (hence completely closed form), while \(K_{-F}\) is the Artin sign Mertens constant in Table \ref{tab:real_input_40_artin_sign_mertens}. Three-factor closed form/numerical verification at constant level one-click generated by script \texttt{scripts/exp\_real\_input\_40\_logM\_artin\_factorization.py}:
\input{sections/generated/eq_real_input_40_logM_artin_factorization_en}
\end{remark}
\begin{corollary}[Certified tail bound for two-kernel series]\label{cor:real-input-40-mertens-two-series-tail}
Let \(q:=z_\star=\varphi^{-2}\).
For any \(N\ge 2\), let
\[
S_3^{(>N)}:=\sum_{m>N}\frac{\mu(m)}{m}\log(1-3q^m+q^{2m}),
\qquad
S_\pm^{(>N)}:=\sum_{m>N}\frac{\mu(m)}{m}\log(1+q^m-q^{2m}).
\]
Then there holds certified upper bound
\[
\boxed{\;
\bigl|S_3^{(>N)}\bigr|+\bigl|S_\pm^{(>N)}\bigr|
\le
5\sum_{m>N}\frac{q^m}{m}
\le
\frac{5}{1-q}\cdot\frac{q^{N+1}}{N+1}\; }.
\]
\end{corollary}
\begin{proof}
For \(t\in[0,q^2]\) direct estimation gives
\(
0\le -\log(1-3t+t^2)\le 4t
\)
and
\(
0\le \log(1+t-t^2)\le t
\).
Taking \(t=q^m\) (when \(m\ge 2\) we have \(t\le q^2\)), and using \(|\mu(m)|\le 1\) gives
\[
\left|\frac{\mu(m)}{m}\log(1-3q^m+q^{2m})\right|\le \frac{4q^m}{m},\qquad
\left|\frac{\mu(m)}{m}\log(1+q^m-q^{2m})\right|\le \frac{q^m}{m}.
\]
Summation yields the first inequality. The second follows from
\(
\sum_{m>N}q^m/m\le \frac{q^{N+1}}{(N+1)(1-q)}
\). This completes the proof.
\end{proof}
\begin{proposition}[Reduced Mertens constant: unique contribution after stripping trivial factors]\label{prop:real-input-40-reduced-mertens}
Suppose
\[
\Delta(z)=\det(I-zM)=\Delta_{\mathrm{triv}}(z)\,\Delta_{\mathrm{nt}}(z),
\qquad
\Delta_{\mathrm{triv}}(z)=(1-z)^a(1+z)^b,
\]
where $a,b\in\ZZ_{\ge 0}$. Assume main pole $z_\star\in(0,1)$ comes from nontrivial block, i.e., $\Delta_{\mathrm{nt}}(z_\star)=0$ is a simple zero, while $\Delta_{\mathrm{triv}}(z_\star)\neq 0$.
Let
\[
\log\mathfrak{M}_{\mathrm{nt}}
:=\log C_{\mathrm{nt}}-\sum_{m\ge 2}\frac{\mu(m)}{m}\log \Delta_{\mathrm{nt}}(z_\star^{\,m}),
\qquad
C_{\mathrm{nt}}:=\frac{1}{-z_\star\,\Delta_{\mathrm{nt}}'(z_\star)}.
\]
Then the Abel finite part constant satisfies strict decomposition
\[
\boxed{\;
\log\mathfrak{M}
\;=\;
\log\mathfrak{M}_{\mathrm{nt}}
\;+\;
(a-b)\,z_\star
\;+\;
b\,z_\star^2.
\;}
\]
In particular, after removing a completely explicit "trivial correction term"\((a-b)z_\star+bz_\star^2\), the real mechanism information of the constant term depends only on nontrivial spectral block $\Delta_{\mathrm{nt}}$.
\end{proposition}
\begin{proof}
From $\Delta=\Delta_{\mathrm{triv}}\Delta_{\mathrm{nt}}$ and $\Delta_{\mathrm{nt}}(z_\star)=0$ being simple zero, we get
\(
\Delta'(z_\star)=\Delta_{\mathrm{triv}}(z_\star)\Delta_{\mathrm{nt}}'(z_\star)
\),
thus main residue constant satisfies
\(
C=\frac{1}{-z_\star\Delta'(z_\star)}=C_{\mathrm{nt}}/\Delta_{\mathrm{triv}}(z_\star)
\).
Substituting
\[
\log\mathfrak{M}=\log C-\sum_{m\ge 2}\frac{\mu(m)}{m}\log\Delta(z_\star^{\,m})
\]
and factoring out $\Delta_{\mathrm{triv}}$ and $\Delta_{\mathrm{nt}}$ gives
\[
\log\mathfrak{M}
=\log\mathfrak{M}_{\mathrm{nt}}
-\log\Delta_{\mathrm{triv}}(z_\star)
-\sum_{m\ge 2}\frac{\mu(m)}{m}\log\Delta_{\mathrm{triv}}(z_\star^{\,m}).
\]
For $\Delta_{\mathrm{triv}}(z)=(1-z)^a(1+z)^b$, using Lemma \ref{lem:mobius-collapse}'s two identities and noting
\(
\sum_{m\ge 2}=\sum_{m\ge 1}-\text{(}m=1\text{ term)}
\),
we can directly simplify to
\[
-\log\Delta_{\mathrm{triv}}(z)-\sum_{m\ge 2}\frac{\mu(m)}{m}\log\Delta_{\mathrm{triv}}(z^{m})
=(a-b)z+bz^2.
\]
Taking $z=z_\star$ gives the conclusion. This completes the proof.
\end{proof}
\begin{definition}[Mertens constant for Artin sign channel]\label{def:real-input-40-artin-sign-mertens}
In tensor decomposition, the factor
\(
\Delta_{\mathrm{twist}}(z)=1+z-z^2
\)
is equivalent to
\(
\det(I-z(-F))=\det(I+zF)
\)
(where \(F=\begin{pmatrix}1&1\\1&0\end{pmatrix}\)).
Denote
\[
\zeta_{-F}(z):=\frac{1}{1+z-z^2},
\qquad
\lambda=\varphi^2,\quad z_\star=\lambda^{-1}.
\]
Define its Abel/Mertens finite part constant (pure twist factor) as
\[
\boxed{\;
K_{-F}
:=\sum_{m\ge 1}\frac{\mu(m)}{m}\log\zeta_{-F}(z_\star^{\,m})
=-\sum_{m\ge 1}\frac{\mu(m)}{m}\log\!\bigl(1+z_\star^{\,m}-z_\star^{\,2m}\bigr)\; }.
\]
\end{definition}
\begin{remark}[Reproducible numerics: \texorpdfstring{$K_{-F}$}{K_-F}]\label{rem:real-input-40-artin-sign-mertens-value}
From Table \ref{tab:real_input_40_finite_part_split}'s splitting, \(\log\mathfrak{M}_{\mathrm{twist}}\) is precisely the above \(K_{-F}\) (it is completely contributed by the non-main-pole factor \(1+z-z^2\)).
For ease of verification, we also provide a high-precision constant table:
\input{sections/generated/tab_real_input_40_artin_sign_mertens_en}
\end{remark}
Additionally, from $\sqrt5=2\varphi-1$, residue constant can also be written as
\[
C=\frac{47+21\sqrt5}{100}=\frac{13+21\varphi}{50}.
\]












